\chapter{Руководство оператора}

Вкладка МКО предназначена для проверки функций микросервиса МКО через графический интерфейс. Левая часть вкладки используется для режима КК, правая часть --- для режима ОУ. В нижней части расположено поле ``Журнал операций'', в котором отображаются действия оператора, результаты вызовов, ошибки gRPC и принятые данные.

\section{Общие правила ввода}

Числовые поля вводятся в соответствии с диапазонами, которые ожидает драйвер и контракт микросервиса. Для слов данных СД допускается ввод как в десятичной, так и в шестнадцатеричной системе счисления. Например, значения \texttt{4660} и \texttt{0x1234} интерпретируются одинаково.

При загрузке слов данных из текстового файла допускаются строки следующих видов: \texttt{0x1234}, \texttt{4660}, \texttt{SD1=0x1234}, \texttt{1 0x1234}. Допускается использовать русское обозначение СД в файле, однако для переносимости рекомендуется использовать латинское обозначение \texttt{SD}.

В строке может быть комментарий после символа \texttt{\#} или последовательности \texttt{//}. Пустые строки игнорируются. Из файла загружается не более 32 слов. Если слов меньше 32, оставшиеся строки таблицы заполняются нулями.

Для двоичных файлов используется формат \texttt{uint32 little-endian}: одно слово СД занимает 4 байта. При импорте из двоичного файла читается не более 32 слов. Если размер файла не кратен 4 байтам, неполный хвост игнорируется, а сообщение об этом выводится в журнал операций.

\section{Диапазоны значений}

\begin{longtable}{|p{0.32\textwidth}|p{0.22\textwidth}|p{0.36\textwidth}|}
	\hline
	\textbf{Поле} & \textbf{Диапазон} & \textbf{Пояснение} \\
	\hline
	\endfirsthead
	\hline
	\textbf{Поле} & \textbf{Диапазон} & \textbf{Пояснение} \\
	\hline
	\endhead
	Номер МКО & 1--2 & Один номер МКО не должен одновременно использоваться в режиме КК и режиме ОУ. Интерфейс автоматически разводит выбранные номера. \\
	\hline
	Канал & 0--1 & 0 --- основной канал, 1 --- резервный канал. \\
	\hline
	\texttt{bus\_control} & 0--1 & Поле управления шиной для настройки КК. \\
	\hline
	Адрес ОУ & 0--30 & Адрес оконечного устройства. \\
	\hline
	Подадрес ОУ & 0--30 & Подадрес области приема или передачи ОУ. \\
	\hline
	Ответное слово, векторное слово, слово самотеста & 0--65535 & Значения являются 16-битными словами. Рекомендуется вводить их в шестнадцатеричном виде. \\
	\hline
	КС1, КС2 & 0--65535 & Командные слова обмена КК. \\
	\hline
	СД1--СД32 & 0--65535 & Слова данных МКО. Интерфейс технически читает \texttt{uint32}, но для корректного слова МКО следует использовать 16-битный диапазон. \\
	\hline
	Формат обмена & 1, 2, 4, 5, 6, 7, 9, 10 & Поддерживаемые драйвером форматы обмена КК. Для каждого формата драйвер дополнительно проверяет допустимое количество слов данных. \\
	\hline
	Remote port & 1--65535 & UDP-порт платы МКО. \\
	\hline
\end{longtable}

\section{Блоки режима КК}

\subsection{ConfigureKK}

Блок ``КК: ConfigureKK'' настраивает плату МКО в режиме контроллера канала. Оператор задает идентификатор платы, номер МКО, канал, параметры управления шиной, адрес ОУ, служебные слова и адрес платы МКО.

Кнопка ``Настроить КК'' выполняет gRPC-вызов \texttt{ConfigureKK}. В журнал операций выводятся параметры вызова и результат: успешное выполнение либо текст ошибки.

\subsection{ConfigureExchange}

Блок ``КК: ConfigureExchange'' задает параметры обмена: формат, командные слова КС1 и КС2, идентификатор операции и массив слов данных СД1--СД32.

Перед таблицей СД расположены кнопки:

\begin{itemize}
	\item ``Загрузить СД из txt'' --- заполняет таблицу из текстового файла;
	\item ``Сохранить СД в txt'' --- сохраняет текущие значения таблицы в текстовый файл;
	\item ``Загрузить СД из bin'' --- заполняет таблицу из двоичного файла;
	\item ``Сохранить СД в bin'' --- сохраняет текущие значения таблицы в двоичный файл.
\end{itemize}

Кнопка ``Настроить обмен'' выполняет gRPC-вызов \texttt{ConfigureExchange}. В журнал выводятся \texttt{board\_id}, формат, КС1, КС2 и результат выполнения.

Кнопка ``RUN: запустить обмен'' выполняет gRPC-вызов \texttt{RunExchange}. Этот вызов отправляет команду запуска обмена в драйвер. В журнал выводятся \texttt{board\_id}, \texttt{operation\_id} и результат выполнения.

Кнопка ``Подписаться на результаты'' фиксирует действие оператора для получения результатов обмена КК. Результаты обмена должны отображаться через поток \texttt{SubscribeExchangeResults} при подключении соответствующей обработки stream-вызова.

\section{Блоки режима ОУ}

\subsection{ConfigureOu}

Блок ``ОУ: ConfigureOu'' настраивает плату в режиме оконечного устройства. Оператор задает идентификатор платы, номер МКО, канал, адрес ОУ, ответное слово и адрес платы.

Кнопка ``Настроить ОУ'' выполняет gRPC-вызов \texttt{ConfigureOu}. В журнал выводятся основные параметры настройки и результат выполнения.

\subsection{SubscribeOuCommands}

Блок ``ОУ: SubscribeOuCommands'' открывает подписку на входящие команды ОУ. Оператор задает \texttt{board\_id} и нажимает ``Подписаться на команды ОУ''. После открытия подписки кнопка меняет назначение на отписку.

При получении события из потока \texttt{SubscribeOuCommands} в журнал выводятся:

\begin{itemize}
	\item командное слово \texttt{cmd\_word} в шестнадцатеричном виде;
	\item слово результата \texttt{result\_word} в шестнадцатеричном виде;
	\item направление обмена;
	\item адрес ОУ;
	\item подадрес;
	\item количество слов;
	\item расшифровка команды;
	\item расшифровка результата.
\end{itemize}

\subsection{ReadOuSubaddress}

Блок ``ОУ: ReadOuSubaddress'' читает выбранный подадрес ОУ. Оператор задает \texttt{board\_id}, номер подадреса и выбирает область: область приема или область передачи.

Кнопка ``Прочитать подадрес'' выполняет gRPC-вызов \texttt{ReadOuSubaddress}. В журнал операций выводится структура ответа: подадрес, командное слово, слово результата, расшифровка результата и массив СД.

\subsection{WriteOuSubaddress}

Блок ``ОУ: WriteOuSubaddress'' записывает слова данных в область передачи выбранного подадреса ОУ. Оператор задает \texttt{board\_id}, подадрес и значения СД1--СД32.

Перед таблицей СД доступны кнопки загрузки и сохранения в текстовом и двоичном формате. Кнопка ``Записать подадрес'' выполняет gRPC-вызов \texttt{WriteOuSubaddress}. В журнал выводятся \texttt{board\_id}, подадрес и переданный массив СД.

\subsection{SendRawOuData}

Блок ``ОУ: SendRawOuData'' предназначен для отправки сырых слов данных ОУ. Поля и файловые операции аналогичны блоку \texttt{WriteOuSubaddress}.

Кнопка ``Отправить сырые данные ОУ'' выполняет gRPC-вызов \texttt{SendRawOuData}. В журнал выводятся \texttt{board\_id}, подадрес и массив переданных слов.

\subsection{SetOuResponseWord}

Блок ``ОУ: SetOuResponseWord'' изменяет ответное слово ОУ. Оператор задает идентификатор платы, номер МКО, канал, адрес ОУ и новое ответное слово.

Кнопка ``Установить ответное слово'' выполняет gRPC-вызов \texttt{SetOuResponseWord}. В журнал выводятся параметры вызова и результат выполнения.

\subsection{ClearReceiveBuffer и ClearTransmitBuffer}

Блоки ``ОУ: ClearReceiveBuffer'' и ``ОУ: ClearTransmitBuffer'' очищают соответственно буфер приема и буфер передачи для выбранной платы и номера МКО.

Кнопки ``Очистить буфер приема'' и ``Очистить буфер передачи'' выполняют gRPC-вызовы \texttt{ClearReceiveBuffer} и \texttt{ClearTransmitBuffer}. В журнал выводятся \texttt{board\_id}, номер МКО и результат операции.

\section{Журнал операций}

Журнал операций используется для контроля действий оператора и диагностики обмена. В журнал выводятся:

\begin{itemize}
	\item факт готовности вкладки и адрес gRPC endpoint;
	\item параметры выполненных вызовов;
	\item сообщения об успешном выполнении команд;
	\item ошибки gRPC и ошибки открытия файлов;
	\item результаты чтения подадреса ОУ;
	\item события входящих команд ОУ из stream-подписки;
	\item сообщения о загрузке и сохранении СД из текстовых и двоичных файлов;
	\item предупреждения о лишних словах в файле или некратном 4 байтам размере двоичного файла;
	\item сообщения об автоматическом изменении номера МКО, если один номер был выбран одновременно для КК и ОУ.
\end{itemize}

Если вызов завершился ошибкой, оператор должен проверить введенные значения, доступность микросервиса МКО, адрес платы и состояние fake-driver или реального драйвера.
